\chapter*{Appendix A: Odoo 19 Community Quick Reference}
\addcontentsline{toc}{chapter}{Appendix A: Odoo 19 Community Quick Reference}
\markboth{Appendix A: Odoo Quick Reference}{}

This appendix condenses the terminology and troubleshooting rules in the supplied West End Bicycles guide.

\section*{Navigation and record behaviour}

\begin{erptable}
\begin{tabularx}{\linewidth}{@{}>{\bfseries\leavevmode\color{ERPNavy}}p{0.22\linewidth}X@{}}
\toprule
\rowcolor{ERPPaleBlue!92}
Term & Meaning \\
\midrule
App grid & Main launcher for Sales, Purchase, Inventory, Manufacturing, CRM, Employees, and other installed apps. \\
Breadcrumb & Clickable trail showing the current app/record path; use it to return to earlier records. \\
Smart button & Button above a record that jumps to related records such as Receipt, Delivery, or Invoices. \\
New & Starts a new record. Existing records are already editable in the Odoo 19 workflow described by the guide. \\
Cloud/save indicator & Appears when changes are unsaved; navigating away also auto-saves in the guide's workflow. \\
Graph / Pivot & Reporting views that summarise transaction data for analysis. \\
\bottomrule
\end{tabularx}
\end{erptable}

\section*{Important Odoo 19 product rule}

The supplied guide states that older ``Storable Product'' terminology should be translated to:

\begin{center}
\textbf{Product Type = Goods} \qquad + \qquad \textbf{Track Inventory = enabled}
\end{center}

If Track Inventory is not enabled, Odoo will not maintain the stock count expected by the course workflow.

\section*{Core transaction vocabulary}

\begin{erptable}
\begin{tabularx}{\linewidth}{@{}>{\bfseries\leavevmode\color{ERPNavy}}p{0.22\linewidth}X@{}}
\toprule
\rowcolor{ERPPaleBlue!92}
Term & Course meaning \\
\midrule
RFQ & Draft purchasing document before supplier commitment. \\
Purchase Order & Confirmed order to a vendor; creates/links a Receipt. \\
Receipt & Inventory transfer bringing purchased goods into stock. \\
BOM & Component recipe for one finished product. \\
Manufacturing Order & Instruction to produce a quantity; consumes components and produces finished stock. \\
Opportunity & Potential sale in CRM. \\
Quotation & Draft commercial offer to a customer. \\
Sales Order & Confirmed customer order; links to fulfilment and billing. \\
Delivery & Inventory transfer sending goods to the customer. \\
Invoice & Customer billing document; Draft before posting, then Posted, then Paid/In Payment after settlement. \\
On Hand & Quantity physically recorded in stock at the current moment. \\
Validate & Complete an inventory transfer and update stock. \\
\bottomrule
\end{tabularx}
\end{erptable}

\section*{Troubleshooting by process state}

\begin{erptable}
\begin{tabularx}{\linewidth}{@{}>{\bfseries\leavevmode\color{ERPRed}}p{0.29\linewidth}p{0.31\linewidth}X@{}}
\toprule
\rowcolor{ERPPaleRed!70}
Problem & Likely cause & What to check \\
\midrule
No Track Inventory option / stock not moving & Product setup or Inventory app issue & Confirm Inventory is installed, Product Type is Goods, and Track Inventory is enabled. \\
No New Quotation from CRM & Opportunity/customer/app state incomplete & Save opportunity, set customer, confirm Sales is installed; otherwise create quotation from Sales. \\
Vendor does not appear on product purchasing & Vendor record exists but is not linked to product & Product Purchase tab, vendor line, price, and Can be Purchased. \\
Cannot receive goods & RFQ not yet confirmed & Confirm RFQ into Purchase Order, then open Receipt. \\
Cannot manufacture & Components are not actually on hand & Validate supplier receipts and check component On Hand. \\
Cannot deliver & Finished product is unavailable & Complete Manufacturing Order and verify Road Bike On Hand. \\
Cannot register payment & Invoice still Draft & Post/confirm the invoice before Register Payment. \\
Expected button missing & Record is in the wrong lifecycle state & Check whether the prior business event is Draft, Confirmed/Posted, or Done. \\
\bottomrule
\end{tabularx}
\end{erptable}

\begin{keyidea}
When troubleshooting, follow the process backwards. A downstream action is usually unavailable because an upstream business event has not been confirmed or completed.
\end{keyidea}
